\chapter{Flujo de agentes}

\section{Cadena cognitiva}

El sistema transforma progresivamente una descripción ambigua en un artefacto
ejecutable. Cada etapa reduce un tipo distinto de incertidumbre.

\begin{figure}[H]
\centering
\begin{verbatim}
Problema natural
  -> paradigmas cientificos
  -> formulaciones matematicas
  -> especificaciones JSON adaptadas
  -> codigo Python + tests
\end{verbatim}
\caption{Cadena cognitiva del pipeline.}
\end{figure}

\section{Clasificación inicial}

Antes de investigar, el sistema intenta asociar el problema a un paraguas
canónico de paradigmas. Esta etapa evita que conceptos equivalentes se almacenen
con nombres distintos. Por ejemplo, una variante de aprendizaje por diferencia
temporal puede pertenecer al paraguas de aprendizaje por refuerzo.

Si el clasificador falla, el sistema degrada de forma segura: continúa con la
investigación sin bloquear la ejecución.

\section{Investigación}

El Researcher combina conocimiento interno y búsqueda externa. Primero consulta
paradigmas conocidos en la memoria. Después usa búsqueda web o académica para
cubrir huecos. Cuando detecta un paradigma relevante, lanza un DeepResearcher
especializado.

\begin{figure}[H]
\centering
\begin{verbatim}
Researcher
  |
  +-- retrieve_knowledge
  +-- web_search / search_papers
  +-- launch_deep_research(paradigma A)
  +-- launch_deep_research(paradigma B)
  |
  v
report.md + deep/{paradigm}.md
\end{verbatim}
\caption{Estructura de la etapa de investigación.}
\end{figure}

El informe profundo incluye fundamentos, postulados, variables, predicciones,
conceptos clave y referencias. El Researcher sintetiza resultados, pero la
información completa queda guardada como artefacto para etapas posteriores.

\section{Formalización}

El Formalizer lanza un subagente por paradigma aprobado. Cada subagente genera
dos o tres formulaciones matemáticas realmente diferentes. No basta con cambiar
nombres: deben existir diferencias estructurales, como una regla de aprendizaje,
una utilidad probabilística o una política umbral.

Cada formulación contiene variables, parámetros, ecuaciones, lógica de decisión
y comparación con alternativas. La salida sigue siendo textual, pero ya está
orientada a implementación.

\section{Especificación del entorno}

El Reasoner necesita conocer el entorno de simulación. Por eso el pipeline
introduce un artefacto \texttt{env\_spec.json} que define acciones disponibles,
percepciones, recursos, restricciones y fuentes de recompensa.

Este archivo funciona como contrato entre ciencia y simulación:

\begin{verbatim}
Formulacion matematica + env_spec.json
        |
        v
Especificacion implementable
\end{verbatim}

\section{Razonamiento de implementación}

El Reasoner adapta cada formulación seleccionada al entorno. Su salida es una
especificación JSON que el Builder puede convertir en código. Además de describir
variables y parámetros, debe mapear cada elemento a percepciones y acciones
observables.

La etapa valida que:

\begin{itemize}
\item las variables usadas estén definidas;
\item las ecuaciones no sean circularmente incoherentes;
\item las acciones existan en el entorno;
\item los valores por defecto sean razonables;
\item la lógica de decisión pueda implementarse.
\end{itemize}

Cuando una formulación no es implementable, el sistema escribe un informe de
validación en lugar de una especificación falsa.

\section{Construcción}

El Builder transforma especificaciones aprobadas en código Python y tests. Cada
subagente escribe un archivo de modelo y un archivo de pruebas. Después ejecuta
las pruebas con \texttt{uv run pytest}.

\begin{figure}[H]
\centering
\begin{verbatim}
spec.json
   |
   v
BuilderSubAgent
   |
   +-- formulation_model.py
   +-- test_formulation.py
   |
   v
uv run pytest
\end{verbatim}
\caption{Construcción y validación ejecutable.}
\end{figure}

La salida aceptada se registra como modelo disponible para el laboratorio. La
fase 2 puede cargarlo porque respeta el contrato \texttt{DecisionModel}.

\section{Revisión humana}

La revisión humana aparece después de las etapas que cambian el significado del
trabajo: investigación, formalización, razonamiento y construcción. El usuario
puede aprobar, rechazar o pedir reejecución. Esta intervención no sustituye a
las pruebas, pero cubre criterios científicos que un test unitario no puede
medir.

\begin{verbatim}
Salida de agente -> revision humana -> memoria -> siguiente etapa
\end{verbatim}

Esta decisión es importante: la memoria persistente representa conocimiento
aceptado, no todos los intentos del modelo.
